\chapter{怎样读出 DNA 序列}

\begin{kaichang}
也许你在新闻里见过这样的东西：一个小小的唾液采集管。你往里面吐一口唾沫，寄回公司，几周后收到一份报告——你的祖先可能来自哪些地区、你对咖啡因的代谢偏快还是偏慢、你携带某些遗传变异的情况。同样的“一口唾沫”戏法还有别的版本：刑侦剧里，一根头发就能锁定嫌疑人；前几年，一根咽拭子就能告诉你有没有感染某种病毒；医院里，一管孕妇的血就能筛查胎儿是否带有某些染色体异常。

请先猜一猜：唾液里明明只有一些脱落细胞和蛋白质，那台神秘的机器是怎么从里面“读”出你 DNA 上的字母的？它真的像扫描仪一样，把 DNA 链一页一页翻过去看吗？

这一章就来拆这台“黑箱”。你会看到，读出 DNA 序列靠的不是“看”，而是一连串巧妙的化学：先\textbf{复制}，再\textbf{分离}，最后让 DNA 在合成过程中\textbf{自己留下线索}。
\end{kaichang}

\section{先把“那一小段”变多：PCR}

第 67 章讲过，细胞复制 DNA 用的是一台叫 DNA 聚合酶的分子机器：它以一条旧链为模板，按照碱基配对规则，一个一个接上核苷酸，造出一条互补的新链。这台机器在试管里也能工作——但有一个麻烦：你的唾液样本里，DNA 总量微乎其微，而真正想研究的往往只是其中一小段，比如某个基因的几百个字母。在大海里捞一根针之前，得先想办法把这根针\textbf{复制出一百亿根}。

这个办法叫聚合酶链式反应，简称 PCR。它的思路朴素得让人意外：既然 DNA 聚合酶只会从某个“起点”开始抄写，那就人为地给它指定起点。

具体做法是这样的。你想复制的那段 DNA，开头和结尾的序列是已知的。于是化学家们合成两小段单链 DNA，每段二十来个字母，分别与目标片段两端的序列互补配对——它们叫\textbf{引物}，相当于贴在目标两端的“从此处开始抄”的标签。把引物、DNA 聚合酶、四种核苷酸原料和样本 DNA 一起放进小试管，然后让温度在三个档位之间循环：

\begin{itemize}[itemsep=2pt]
\item 加热到 $94\,^{\circ}\mathrm{C}$ 左右：双链 DNA 的氢键被热运动拆散，两条链分开——这一步叫变性；
\item 降温到 $55\,^{\circ}\mathrm{C}$ 左右：引物恰好能稳稳地贴在各自的目标位置上——这一步叫退火；
\item 升到 $72\,^{\circ}\mathrm{C}$ 左右：聚合酶从引物末端开始，沿模板链合成新链——这一步叫延伸。
\end{itemize}

一轮循环下来，原来的 1 份目标片段变成了 2 份。再来一轮，2 变 4，4 变 8……每一轮，上一轮造出来的新链也会充当模板。引物保证了：被疯狂复制的只有夹在两个引物之间的那一段，基因组里其余几十亿个字母都被晾在一边。这就是 PCR “选择性”的来源——\textbf{选择性不靠机器聪明，靠的是碱基配对规则本身}：引物只会贴在序列互补的地方。

\begin{qiaoliang}
PCR 的扩增是指数式的：$n$ 轮循环之后，目标片段的理论份数是 $2^n$。算一笔账：$2^{10}\approx 10^3$（一千），$2^{20}\approx 10^{6}$（一百万），$2^{30}\approx 10^{9}$（十亿）。也就是说，哪怕样本里最初只有几十个目标分子，30 轮循环——在自动仪器上大约一两个小时——之后就有几十亿份，多到能被检测出来。这正是核酸检测能“从一口痰里找到一个病毒”的底气。

当然，真实的 PCR 达不到完美的翻倍：酶不是每次都成功，原料会逐渐耗尽，产物太多还会互相干扰。如果每轮的平均扩增倍数是 1.8 而不是 2，那么 30 轮后是 $1.8^{30}\approx 4.6\times 10^{7}$——只有理论值的大约二十分之一，但仍然足够用。指数增长的威力在于：哪怕底数小一点，多跑几轮就能补回来；反过来，最初的目标分子差 10 倍，最终产量也恰好差 10 倍——所以 PCR 还能用来\textbf{数}样本里原来有多少目标分子（定量 PCR），比如估计病毒感染的严重程度。
\end{qiaoliang}

这里有一个很容易被忽略的问题：$94\,^{\circ}\mathrm{C}$ 会把绝大多数蛋白质“煮熟”——第 48、49 章讲过，蛋白质的功能依赖折叠形状，高温让酶变性失活。你身体里的 DNA 聚合酶在这个温度下第一轮就报废了。PCR 能成立，要感谢一种生活在温泉里的细菌——水生栖热菌，它在美国黄石公园近沸腾的热泉里过得很好，它的 DNA 聚合酶（简称 Taq 酶）天生耐热，反复经受 $94\,^{\circ}\mathrm{C}$ 也不罢工。没有它，每一轮循环都得重新加酶，PCR 就只是论文里的一个想法，而不会变成今天全世界实验室的日常工具。这也是第 76、77 章讲进化时埋下的一个注脚：生命在极端环境里演化出的“绝活”，常常成为人类技术的宝库。

PCR 的应用几乎无处不在：刑侦现场一根头发里的一点 DNA，扩增后足以比对身份；古生物学家从几万年前的尼安德特人骨头里扩增出残存的 DNA 片段；医生用它查病原体；生态学家从一勺湖水或一撮土里扩增 DNA，盘点里面住着哪些看不见的生命（这叫环境 DNA 技术）。但请记住：PCR 只负责\textbf{复制}，它本身并不能告诉你序列是什么。读出字母，还需要后面几步。

\section{按个头排队：凝胶电泳}

下一步要解决的是分离。想象你手里有一堆长短不一的 DNA 片段，怎么知道它们各有多长？答案是让它们\textbf{赛跑}。

先把镜头推到粒子层。DNA 的骨架由糖和磷酸交替连接而成（第 66 章），磷酸基团在细胞的酸碱环境下会解离出 \chem{H^+}，使骨架上布满负电荷。也就是说，每一条 DNA 链都是一个长长的负离子。把它放进电场里，它就会朝正极移动——电荷越多的片段受力越大，而片段越长磷酸越多，负电荷也越多，粗看下来大家的“推动力”似乎差不多。

真正的筛选靠的是跑道。把 DNA 加进一块琼脂糖凝胶的一端——琼脂糖是从海藻里提取的多糖，煮成溶液再冷却，会形成一张充满纳米级孔洞的三维网，像一块极其细密的海绵。接通电场后，DNA 负离子在这张网里钻向正极：短片段个头小，在孔洞里钻得快；长片段不断被网挂住，走得慢。跑上一段时间，不同长度的 DNA 就按个头排成了一道道“条带”：最前面是最短的，最后面是最长的。旁边通常会加一列已知长度的标准片段作为“梯子”，一对照，就能读出每条带的长度。

\begin{center}
\begin{tikzpicture}[scale=0.9]
\draw[fill=gray!15] (0,0) rectangle (7,4);
\foreach \x in {1,2.5,4,5.5} {\draw[fill=white] (\x-0.35,3.6) rectangle (\x+0.35,3.9);}
\draw[thick] (0.65,2.9) -- (1.35,2.9);
\draw[thick] (0.65,2.2) -- (1.35,2.2);
\draw[thick] (0.65,1.5) -- (1.35,1.5);
\draw[thick] (2.15,2.2) -- (2.85,2.2);
\draw[thick] (3.65,2.9) -- (4.35,2.9);
\draw[thick] (3.65,1.5) -- (4.35,1.5);
\draw[thick] (5.15,2.9) -- (5.85,2.9);
\draw[thick] (5.15,2.2) -- (5.85,2.2);
\draw[thick] (5.15,1.5) -- (5.85,1.5);
\draw[thick] (5.15,0.8) -- (5.85,0.8);
\node at (3.5,-0.5) {正极方向（下方），条带越靠下，DNA 片段越短};
\end{tikzpicture}
\end{center}

这是一个人为的表示法：真实的凝胶是半透明的，条带要靠能与 DNA 结合的染料染色后，在特定光线下才看得见；DNA 分子本身也不是小横线，而是卷曲的长链。电泳装置使用不低的电压，染料中有些品种对健康有害，所以\textbf{电泳只能在实验室由老师或专业人员操作}，请不要在家尝试。不过，有一个和它开头几步相同的观察，你可以安全地亲手做：

\begin{shiyan}
\textbf{从草莓里“捞出”DNA（家庭安全版）}

材料：熟透的草莓 2～3 颗、食盐一小撮、洗洁精一小勺、清水、自封袋、纱布或咖啡滤纸、透明玻璃杯、冰镇的高度酒精（如冰镇医用酒精，需家长协助，全程远离火源）。

步骤：①把草莓装进自封袋，挤出空气封好，用手捏成糊状——这一步是物理破碎细胞；②小半杯温水里溶入盐和洗洁精，倒入袋中混匀，静置 10 分钟——洗洁精破坏细胞膜和核膜的脂质双层（回忆第 47 章），盐帮助 DNA 聚集；③把袋中液体隔着纱布过滤到玻璃杯里；④沿杯壁缓缓倒入等体积的冰酒精，不要搅拌，静置几分钟。

观察点：在果汁与酒精的交界处，会出现一团团白色、半透明的絮状物，用牙签可以挑起来——那就是成千上万条缠在一起的 DNA（混着一些蛋白质）。DNA 不溶于冰冷的酒精，于是从水里析出。

安全提示：酒精易燃，全程远离明火；实验后的材料都不能吃；洗洁精不要入眼。想一想：我们看到的絮状物是“一条 DNA 分子”吗？（提示：一条 DNA 分子的直径只有 2 纳米。）
\end{shiyan}

凝胶电泳能做的远不止量长度：PCR 做完后先跑一块胶，确认想要的片段确实被扩增出来了（只有目标长度处出现条带，说明引物设计得靠谱）；比较不同个体在同一段 DNA 上的差异；检查 DNA 样本是否完整。它至今仍是分子生物学实验室每天必做的基本功。而它最深远的用武之地，是接下来要讲的——读出序列本身。

\section{桑格测序：让 DNA 在“断头路”上自首}

现在面对核心难题：一段几百个字母的 DNA，怎样知道每个位置上是 A、T、G 还是 C？没有任何显微镜能做到这件事——碱基之间的差别只有几个原子，远小于光学显微镜的分辨极限，而电子显微镜也没法可靠地区分四种碱基。直接的“看”走不通，只能让 DNA \textbf{间接地招供}。

1977 年，英国生物化学家弗雷德里克·桑格想出了一个绝妙的主意：\textbf{双脱氧链终止法}，后来干脆就叫桑格测序。回忆第 67 章：DNA 聚合酶接新核苷酸时，是把新核苷酸接到上一条链末端的羟基上，像接龙一样一节一节延长。桑格在试管里除了正常的四种核苷酸，还掺入少量“动过手脚”的核苷酸——它们正好缺了那个供下一个核苷酸连接的羟基。这种“双脱氧核苷酸”一旦被聚合酶接进链里，链条就到了\textbf{断头路}，再也无法延长，反应戛然而止。

关键设计在于：给四种双脱氧核苷酸分别染上四种不同颜色的荧光标记——A 一种颜色，T 一种颜色，G、C 各一种。然后把它们和正常核苷酸、聚合酶、引物、模板混在一起。由于“动手脚”的分子只占少数，每条新链在哪个位置撞上终止核苷酸完全是碰运气：有的在复制到第 50 个字母时碰上了一个终止 A，断在 50；有的在第 51 位碰上终止 G，断在 51；有的直到第 300 位才碰上终止 C。反应结束后，试管里就有了\textbf{各种长度的新链}，而每条链的末尾颜色和它的“死因”一一对应：末尾是红色的，说明它断在一个 A 上。

接下来，把这一大堆断链送去跑“电泳”——实际仪器用的是极细的石英毛细管，原理和凝胶电泳一模一样：按长度排队。最短的片段（只延长了 1 个字母就断掉的）最先到达终点，被激光逐个照亮，仪器读出它的颜色；然后是长 2 个字母的、长 3 个字母的……于是你得到一串颜色序列，而颜色序列翻译过来就是字母序列：第 1 个到达的是链上第 1 个字母，第 2 个到达的是第 2 个字母，依此类推。

\begin{zhengju}
桑格测序的故事值得细讲，因为它不是“天才灵光一闪”，而是一场“怎样读出看不见的信息”的攻关战。

桑格这个人此前已经干成过一件类似的事：1950 年代，他花了十年时间，测出了胰岛素这种蛋白质里全部 51 个氨基酸的排列顺序——第一次证明蛋白质分子有确定不变的序列，为此获得 1958 年诺贝尔化学奖。蛋白质能测，DNA 为什么不能？问题在于 DNA 的字母只有 4 种，链条却长达成千上万，蛋白质测序那套“逐段切下来鉴定”的办法根本忙不过来。到 1970 年代，全世界好几个实验室都在攻关，主流思路是“把 DNA 打断、分别测、再拼”，进展缓慢。

桑格换了一个角度：与其把长链拆了读，不如\textbf{让复制过程本身当报告员}。他先发明了“加减法”，很快又改进为链终止法——利用聚合酶忠实但会“上当”的脾气，把序列信息翻译成“断链长度加末端颜色”，而后者是当时就能测的东西。1977 年论文发表。几乎同一时期，美国的马克萨姆和吉尔伯特发明了另一条路线：用化学试剂有选择地在特定碱基处把 DNA 直接切断，同样靠电泳读长度。两种方法一度并行竞争——这正是健康的科学：同一个问题，两种独立方案互相校验。几年之后胜负分明：化学切割法要用剧毒试剂、步骤繁琐，而链终止法安全、简单、可以自动化。1980 年，桑格因此第二次获得诺贝尔化学奖。

故事还有下半场。1990 年启动的人类基因组计划，要把这套方法用到极致：测出人类全部约 30 亿个字母。公共计划的策略是“先把基因组分成带路标的大片段，逐个测再按路标拼接”，稳妥但慢；1998 年，企业家文特尔创办公司，主张“鸟枪法”——把整个基因组随机打碎成小段全部测掉，交给计算机拼。两派公开争论鸟枪法能不能搞定充满重复序列的人类基因组，这场竞争反而把进度大大提前：2001 年双方同时公布了草图。事后的分析表明，公共计划的“路标”数据对拼通重复区域至关重要——争论双方的方法，最终合流成了今天测序的标准做法。
\end{zhengju}

桑格测序一次能读出约 500～1000 个字母，准确率很高，至今仍是验证小片段序列的“金标准”。但它一次只读一条链。要测一个人的整个基因组——30 亿个字母——靠它一条一条读，需要难以想象的时间和金钱。人类基因组计划当年花了约 13 年、近 30 亿美元。读出更多生命，需要一次质的飞跃。

\section{高通量测序：一百万条链同时开工}

质的飞跃来自一个朴素的想法：\textbf{并行}。既然读一条链的化学已经成熟，那就让一百万条、一亿条链同时读。

以目前应用最广的“边合成边测序”为例。先把 DNA 样本打断成一两百个字母的短片段，把它们密密麻麻地固定在一块玻璃芯片的表面，每个位置一簇，就像芯片上长满了微小的“草丛”。然后在芯片上流过反应液，让每一簇 DNA 同步进行合成反应——但这次掺入的核苷酸带着荧光标记，而且每接上一个，反应就被化学“快门”暂停一次。每暂停一次，仪器就给整块芯片拍一张荧光照片：哪个位置亮了什么颜色，就说明那一簇在这一轮接上了哪种碱基。洗掉标记，恢复合成，再停，再拍照……循环一百多次，一张照片摞一张照片，每一簇的序列就从颜色变化里被读了出来。

一块芯片上有几亿簇，一次运行就产生几亿条序列——尽管每条只有一百多个字母。这就是为什么它叫“高通量”：单条读得短，但总量惊人。对比一下：

\begin{table}[htbp]
\centering
\begin{tabular}{lll}
\toprule
 & 桑格测序 & 高通量测序 \\
\midrule
每次读几条链 & 1 条 & 数千万至数亿条 \\
每条读多长 & 500～1000 个字母 & 100～300 个字母 \\
擅长的事 & 精确验证一小段 & 海量覆盖整个基因组 \\
典型用途 & 确认某个突变 & 基因组、微生物群落 \\
\bottomrule
\end{tabular}
\end{table}

效果有多惊人？2001 年测一个人类基因组要花上亿美元；到 2020 年代，商业价格已降到几百美元，一天就能完成——二十年间成本下降的速度，比著名的计算机芯片“摩尔定律”还快得多。测序就这样从国家级工程变成了医院的常规检查：医生给肿瘤组织测序寻找驱动突变，疾控人员给病毒测序追踪传染病的传播链条，生态学家给一管海水测序清点其中的生命。

\section{把撕碎的书拼回去：组装、错误与偏差}

高通量测序给你的不是一本书，而是几亿张从书上撕下来的小纸条，每张一百来个字。怎么把它们还原成书？这一步叫\textbf{基因组组装}，它本质上是一个巨型拼图游戏：如果纸条甲的末尾 30 个字母和纸条乙的开头 30 个字母一模一样，它们很可能来自书上相邻的位置，可以搭接起来；如此不断重叠、延伸，小段接成大段，大段接成染色体。

这个游戏有两个可怕的难点。第一是\textbf{重复}：人类基因组里超过一半的序列属于各种重复出现的段落，有些一模一样的句子在书里出现几千次。一本侦探小说里每章都出现“那天晚上下着雨”，你手里又只有一百个字的小纸条——你根本分不清这张纸条来自哪一章。解决的办法是“留线索”：测序时有意记录每对纸条来自同一根较长片段的两端（双端测序），撕碎时保留碎片之间的距离信息；近年兴起的“长读长”测序技术则干脆把纸条加长到上万个字母，让大多数重复段落能被一次性跨过。第二个难点是\textbf{错误}：每一张纸条上的字母都有约百分之一的概率读错——荧光拍照、化学反应都不完美。单个字母错了怎么办？

\begin{qiaoliang}
答案是\textbf{多读几遍，举手表决}。假设测序仪每次读一个字母的错误率是 $1\%$，只读 1 遍，这个位置就有 $1\%$ 的可能出错。但如果同一个位置被 30 张独立的纸条覆盖（这叫 30 倍覆盖度），而它们的错误互不相干，那么“超过一半纸条同时读错”的概率低到可以忽略——粗略地说，要让 30 张里至少 16 张同时犯同样的错，概率不到万亿分之一量级。所以临床上测序人类基因组通常要求 30 倍以上的覆盖度：不是仪器不够准，而是用\textbf{冗余}把随机错误投票投出去。反过来，覆盖度低的区域（比如只有 3、4 张纸条），一个测序错误就可能被当成真实突变——这是解读测序数据时最经典的陷阱之一。
\end{qiaoliang}

组装还有一个更隐蔽的问题：\textbf{参考基因组偏差}。实际操作中，为了省事，人们常常不做完整的拼图，而是把纸条直接贴到一本“标准答案”——参考基因组——上，只记录和标准答案不同的地方。这很快，但代价是：凡是你的 DNA 里有而参考基因组里没有的大段内容，纸条根本找不到贴的位置，就被丢弃了。更值得警惕的是，第一版人类参考基因组主要来自少数几位志愿者，其中约七成来自同一个人。拿它当“人类的标准”，本身就带着偏差：不同人群特有的序列变异在参考基因组里缺席，相应人群的数据就更容易被误读。科学界正在用成千上万人的基因组拼凑更能代表人类多样性的“泛基因组”参考——承认“标准答案”本身也在被修正，这正是科学诚实的样子。

\section{边界：读出字母，不等于读懂意思}

到这里，黑箱基本拆完了：PCR 负责放大目标，电泳负责按大小分离，链终止和荧光负责让 DNA 自首，计算机负责拼图。这套技术何时可靠、何时会失效，值得认真清点：

\begin{itemize}[itemsep=2pt]
\item 引物设计不好，PCR 会扩增出错误的片段，或者干脆什么都扩不出来——所以电泳验证不可省略；
\item 高重复序列的区域，短读长纸条拼不进去，组装出的基因组总有缺口；
\item 样本被污染（比如混入了细菌的 DNA），测序仪会忠实地把污染也读出来；
\item 覆盖度不足时，测序错误会伪装成“突变”；
\item 对着参考基因组比对，会漏掉个体独有的大段结构差异。
\end{itemize}

\begin{wuqu}
“现在技术这么发达，测出一个人的全部 DNA 序列，就等于拿到了这个人生命的说明书，能预言他会得什么病、成为什么样的人。”——这是对测序最常见的夸大。第一，序列本身都未必读全读对（覆盖度、重复区、参考偏差，前文刚讲过）。第二，就算字母全对，\textbf{字母的含义大半还是谜}：人类基因组里只有约百分之一出头直接编码蛋白质，其余大量序列的功能至今不清楚；同一个字母差异，在一个人身上有显著影响，在另一个人身上可能毫无表现——因为性状是基因、基因之间、基因与环境共同编织的网（第 73 章）。第三，序列之外还有一整层“使用方式”的信息：同样的序列，哪些基因被打开、开到多大，取决于表观遗传标记（第 74 章），而测序读不出这些标记。DNA 序列更像一本只译出了百分之一的古书——拿到原文，离读懂还远。
\end{wuqu}

\section{回到那管唾液}

现在可以回答开场的问题了。你的唾液里有从口腔黏膜脱落的细胞，每个细胞的细胞核里都有一份完整的基因组。检测公司提取出 DNA，用 PCR 把报告要分析的那些片段成亿倍地放大（或者用高通量测序整体读一遍），读出特定位点上的字母，再和数据库里大量人群的数据比较，推算出“这个位点的这种版本，在某人群中更常见”“有这种版本的人，咖啡因代谢平均偏慢”之类的结论。

注意其中的每一步都建立在概率和统计上：扩增有效率，测序有错误，比对有偏差，而最后的“解读”更是来自人群层面的相关性，不是关于你个人的铁律。那份报告可以是有趣的提示，但它远没有“看穿”你——也不可能看穿你。

这条链还没有走完。读，是为了理解；而理解之后，人类很快就开始问一个更大胆的问题：既然能读，能不能\textbf{写}？能不能精确地改掉一个字母、一段基因？这正是下一章的主题——改写基因的技术，以及它带来的希望与缰绳。

\begin{tiaozhan}
为什么短读长拼不过重复区？可以算一笔粗账。设基因组某段有一个长度为 $L$ 的重复序列，在基因组里出现了两次且完全相同。如果测序纸条的长度 $r<L$，那么完全落在重复区内部的纸条，在书上有两个无法区分的出处，拼接程序只能停下。只有当纸条长度超过重复区，长到两端都伸进了重复区两侧独一无二的“单拷贝”序列里，它的位置才被唯一锁定。人类基因组里存在长达几十万字母、相似度超过 $99\%$ 的重复段落，这就是为什么一百多字母的短读长无论怎样加覆盖度也填不平所有缺口——缺的不是数量，而是单张纸条的“视野”。长读长测序把 $r$ 提高两个数量级，直接绕开了这个数学瓶颈。跳过本段不影响理解主线，但它说明：测序的极限，很多时候是组合学问题，而不只是化学问题。
\end{tiaozhan}

\section*{练一练}

\begin{sikaoti}
\item 画一张 PCR 三轮循环的信息流图：用两条线代表双链 DNA 模板，标出引物的位置和方向，数出第一轮、第二轮、第三轮之后各有几份目标片段。图旁注明：这是表示法，真实的 DNA 不是直线。
\item 在凝胶电泳中，为什么带负电的 DNA 片段越长反而跑得越慢？请用“推动力”和“网的阻力”两股因素分别说明：哪一股随长度怎么变，最终谁占了上风？
\item 一段目标 DNA 最初只有 100 个分子。假设 PCR 每轮扩增 1.8 倍，估算 20 轮后大约有多少个分子？如果最初只有 10 个分子呢？这个结果说明 PCR 的产物数量保留了关于初始量的什么信息？
\item 给你三张纸条：甲 = ……GATTACA，乙 = TACAGCC……，丙 = ……CAGCCAT。请根据重叠把它们拼成一条更长的序列，并写出拼接结果。再回答：如果基因组里“TACA”出现了 50 次，你的拼接还可靠吗？为什么？
\item 某基因检测报告称：“你携带某基因的某版本，某疾病风险是平均水平的 1.5 倍。”结合本章和第 73 章的内容，列出至少三个在据此下结论之前必须追问的问题（提示：覆盖度？相关性还是因果？群体还是个人？）。
\end{sikaoti}
